\chapter{EKLER}

\section*{Ek A. Bağlantılar}
\addcontentsline{toc}{section}{Ek A. Bağlantılar}

{\small\begin{tabularx}{\textwidth}{@{}l X@{}}
Kaynak kodu (GitHub) & \url{https://github.com/AhmedElghazzawy/ai-ticket-assistant} \\
Canlı demo & \url{https://bseu-helpdesk.onrender.com} \\
Personel demosu & \url{https://bseu-helpdesk.onrender.com/staff} $\rightarrow$ ``Demo olarak dene'' \\
Mimari belgesi & \texttt{ARCHITECTURE.md} \\
Karar kayıtları & \texttt{docs/decisions/} (18 ADR) \\
Değerlendirme sonuçları & \texttt{eval/} \\
\end{tabularx}}

Ücretsiz sunucu bir süre ziyaretçi olmazsa uyuduğu için ilk açılış yaklaşık bir dakika sürebilir.

\section*{Ek B. Karar Fonksiyonu}
\addcontentsline{toc}{section}{Ek B. Karar Fonksiyonu}

Bölüm~\ref{sec:karar}'de anlatılan yedi kontrolün koddaki hâli Kod~\ref{lst:karar}'de verilmiştir.

\begin{lstlisting}[language=Python, caption={Otomatik gönderim kararı (\texttt{domain/decision.py})}, label={lst:karar}]
MIN_SCORE = 0.65

def decide(result: Classification, top_score: float, answer: Answer,
           retrieved: dict[str, str]) -> Decision:
    """Auto-send only if ALL checks pass. Every failed check is recorded as a reason."""
    reasons: list[SimpleReason | WeakMatch] = []
    if result.category == "other":
        reasons.append(SimpleReason(code="category_other"))
    if result.escalate:
        reasons.append(SimpleReason(code="escalation"))
    if result.priority == "urgent":  # a person can act today
        reasons.append(SimpleReason(code="urgent"))
    if top_score < MIN_SCORE:
        reasons.append(WeakMatch(score=top_score, min=MIN_SCORE))
    if not answer.covered:
        reasons.append(SimpleReason(code="not_covered"))
    if not answer.sources or not set(answer.sources) <= retrieved.keys():
        reasons.append(SimpleReason(code="no_citation"))
    elif not is_grounded(answer.evidence, list(retrieved.values())):  # e.g. a fake rule
        reasons.append(SimpleReason(code="ungrounded"))
    return Decision(auto_send=not reasons, reasons=reasons)
\end{lstlisting}

\section*{Ek C. Kaynaklı Cevabın Çıktı Şeması}
\addcontentsline{toc}{section}{Ek C. Kaynaklı Cevabın Çıktı Şeması}

\begin{lstlisting}[language=Python, caption={Cevap adımının çıktı şeması (\texttt{domain/outputs.py}, alan açıklamaları Türkçe yorum olarak kısaltılmıştır)}, label={lst:cevap}]
class Answer(BaseModel):
    """`evidence` comes first: quote the rule, then decide, then write."""
    model_config = ConfigDict(extra="forbid", strict=True)

    evidence: str      # kuraldan kelimesi kelimesine alıntı (yoksa boş)
    covered: bool      # belgeler ana soruyu cevaplıyor mu?
    reply: str         # yalnızca belgelerdeki bilgilerle kısa, nazik cevap
    sources: list[str] # kullanılan parça kimlikleri, ör. "sinavlar.tr#4"
    general_help: str  # belgelerin cevaplamadığı kısım için genel tavsiye
\end{lstlisting}

\section*{Ek D. Kararın Testle Güvence Altına Alınması}
\addcontentsline{toc}{section}{Ek D. Kararın Testle Güvence Altına Alınması}

Her kural için önce başarısız olan bir test yazılmıştır. Kod~\ref{lst:test}, talebe yerleştirilmiş
sahte bir kuralı alıntılayan cevabın gönderilmediğini ve öğrenciye ön bilgi olarak da
gösterilmediğini doğrulayan gerçek testi göstermektedir (\texttt{tests/test\_decision.py}).

\begin{lstlisting}[language=Python, caption={Birim testi: belgede bulunmayan alıntı otomatik gönderilmez}, label={lst:test}]
RULE = "No separate application is needed for the make-up exam."
RETRIEVED = {"sinavlar.en#4": f"Make-up exam\n{RULE}", "sinavlar.en#1": "Exam dates are on OBS."}

def test_evidence_not_found_in_the_rules_is_not_auto_sent() -> None:
    # Prompt injection: the ticket carried a fake rule and the model quoted it, citing a real id.
    fake = Answer(
        evidence="Every student who withdraws gets a full refund.",
        covered=True,
        reply="You will get a full refund.",
        sources=["sinavlar.en#4"],
        general_help="",
    )
    decision = decide(labels(), 0.80, fake, RETRIEVED)
    assert not decision.auto_send
    assert [reason.code for reason in decision.reasons] == ["ungrounded"]
    assert not show_first_help(decision)  # the fake text must not reach the student either
\end{lstlisting}
